\documentclass[10pt, letterpaper]{article}

% Packages:
\usepackage[
    ignoreheadfoot, % set margins without considering header and footer
    top=0.5 cm, % seperation between body and page edge from the top
    bottom=1.3 cm, % seperation between body and page edge from the bottom
    left=2 cm, % seperation between body and page edge from the left
    right=2 cm, % seperation between body and page edge from the right
    footskip=1.0 cm, % seperation between body and footer
    % showframe % for debugging 
]{geometry} % for adjusting page geometry
\usepackage{titlesec} % for customizing section titles
\usepackage{tabularx} % for making tables with fixed width columns
\usepackage{array} % tabularx requires this
\usepackage[dvipsnames]{xcolor} % for coloring text
\definecolor{primaryColor}{RGB}{0, 0, 0} % define primary color
\usepackage{enumitem} % for customizing lists
\usepackage{fontawesome5} % for using icons
\usepackage{amsmath} % for math
\usepackage[
    pdftitle={Moad Tellami},
    pdfauthor={Moad Tellami},
    pdfcreator={LaTeX with RenderCV},
    colorlinks=true,
    urlcolor=primaryColor
]{hyperref} % for links, metadata and bookmarks
\usepackage[pscoord]{eso-pic} % for floating text on the page
\usepackage{calc} % for calculating lengths
\usepackage{bookmark} % for bookmarks
\usepackage{lastpage} % for getting the total number of pages
\usepackage{changepage} % for one column entries (adjustwidth environment)
\usepackage{paracol} % for two and three column entries
\usepackage{ifthen} % for conditional statements
\usepackage{needspace} % for avoiding page brake right after the section title
\usepackage{iftex} % check if engine is pdflatex, xetex or luatex

% Ensure that generate pdf is machine readable/ATS parsable:
\ifPDFTeX
    \input{glyphtounicode}
    \pdfgentounicode=1
    \usepackage[T1]{fontenc}
    \usepackage[utf8]{inputenc}
    \usepackage{lmodern}
\fi

\usepackage{charter}

% Some settings:
\raggedright
\AtBeginEnvironment{adjustwidth}{\partopsep0pt} % remove space before adjustwidth environment
\pagestyle{empty} % no header or footer
\setcounter{secnumdepth}{0} % no section numbering
\setlength{\parindent}{0pt} % no indentation
\setlength{\topskip}{0pt} % no top skip
\setlength{\columnsep}{0.15cm} % set column seperation
\pagenumbering{gobble} % no page numbering

\titleformat{\section}{\needspace{3\baselineskip}\bfseries\large}{}{0pt}{}[\vspace{1pt}\titlerule]

\titlespacing{\section}{
    % left space:
    -1pt
}{
    % top space:
    0.15 cm
}{
    % bottom space:
    0.15 cm
} % section title spacing

\renewcommand\labelitemi{$\vcenter{\hbox{\small$\bullet$}}$} % custom bullet points
\newenvironment{highlights}{
    \begin{itemize}[
        topsep=0.10 cm,
        parsep=0.10 cm,
        partopsep=0pt,
        itemsep=0pt,
        leftmargin=0 cm + 10pt
    ]
}{
    \end{itemize}
} % new environment for highlights


\newenvironment{highlightsforbulletentries}{
    \begin{itemize}[
        topsep=0.10 cm,
        parsep=0.10 cm,
        partopsep=0pt,
        itemsep=0pt,
        leftmargin=10pt
    ]
}{
    \end{itemize}
} % new environment for highlights for bullet entries

\newenvironment{onecolentry}{
    \begin{adjustwidth}{
        0 cm + 0.00001 cm
    }{
        0 cm + 0.00001 cm
    }
}{
    \end{adjustwidth}
} % new environment for one column entries

\newenvironment{twocolentry}[2][]{
    \onecolentry
    \def\secondColumn{#2}
    \setcolumnwidth{\fill, 4.5 cm}
    \begin{paracol}{2}
}{
    \switchcolumn \raggedleft \secondColumn
    \end{paracol}
    \endonecolentry
} % new environment for two column entries

\newenvironment{threecolentry}[3][]{
    \onecolentry
    \def\thirdColumn{#3}
    \setcolumnwidth{, \fill, 4.5 cm}
    \begin{paracol}{3}
    {\raggedright #2} \switchcolumn
}{
    \switchcolumn \raggedleft \thirdColumn
    \end{paracol}
    \endonecolentry
} % new environment for three column entries

\newenvironment{header}{
    \setlength{\topsep}{0pt}\par\kern\topsep\centering\linespread{1.5}
}{
    \par\kern\topsep
} % new environment for the header

\newcommand{\placelastupdatedtext}{% \placetextbox{<horizontal pos>}{<vertical pos>}{<stuff>}
  \AddToShipoutPictureFG*{% Add <stuff> to current page foreground
    \put(
        \LenToUnit{\paperwidth-2 cm-0 cm+0.05cm},
        \LenToUnit{\paperheight-1.0 cm}
    ){\vtop{{\null}\makebox[0pt][c]{
        \small\color{gray}\textit{Dernière mise à jour en juillet 2026}\hspace{\widthof{Dernière mise à jour en juillet 2026}}
    }}}%
  }%
}%

% save the original href command in a new command:
\let\hrefWithoutArrow\href


% =========================== DOCUMENT ===============================================
\begin{document}
    \newcommand{\AND}{\unskip
        \cleaders\copy\ANDbox\hskip\wd\ANDbox
        \ignorespaces
    }
    \newsavebox\ANDbox
    \sbox\ANDbox{$|$}

    % =========================== HEADER ==========================
    \begin{header}
        \fontsize{25 pt}{25 pt}\selectfont Moad Tellami

        \vspace{5 pt}
        \normalsize
        \mbox{\hrefWithoutArrow{mailto:moadtell@gmail.com}{moadtell@gmail.com}}%
        \kern 5.0 pt%
        \AND%
        \kern 5.0 pt%
        \mbox{\hrefWithoutArrow{tel:+212693566907}{+212 693 56 69 07}}%
        \kern 5.0 pt%
        \AND%
        \kern 5.0 pt%
        \mbox{\hrefWithoutArrow{https://www.linkedin.com/in/mtellami}{linkedin/mtellami}}%
    \end{header}
    \vspace{8 pt}

    % =========================== SUMMARY ==========================
     Ingénieur logiciel avec plus de 2 ans d'expérience dans le développement d'applications full-stack utilisant Django, NestJS, Angular, React et PostgreSQL. Expérimenté dans les systèmes géospatiaux, l'intégration de l'IA, l'authentification et les services backend évolutifs, avec un accent sur la conception de logiciels maintenables et la résolution de problèmes d'ingénierie complexes. \vspace{5 pt}


    % =========================== EXPERIENCE ==========================
    \section{Expérience}
        \begin{twocolentry}{
            Jan 2026 – Juil 2026
        }
        \textbf{Stagiaire Ingénieur Logiciel AI Delivery} -- \textbf{\hrefWithoutArrow{https://www.bcg.com/x}{BCG}} \end{twocolentry}
        \vspace{0.10 cm}

        \begin{onecolentry}
            \begin{highlights}
                \item Développé des visualisations \textbf{géospatiales} interactives pour des jeux de données environnementaux, notamment les feux de forêt, les tremblements de terre, les sécheresses et d'autres données de risques, permettant aux utilisateurs d'explorer, de comparer et d'analyser des informations de risque en couches.
                \item Intégré des modèles d'\textbf{IA} dans les flux de production pour automatiser le traitement et la classification des données, exposant les prédictions des modèles via le pipeline d'analyse de l'application.
                \item Collaboré étroitement avec les designers \textbf{UI/UX} pour implémenter des interfaces \textbf{React} et \textbf{Tailwind CSS} réactives, traduisant les spécifications de conception en expériences utilisateur soignées et conformes aux exigences des clients.
                \item Implémenté des endpoints \textbf{REST API} avec \textbf{Django} REST Framework et \textbf{PostgreSQL}, en suivant les bonnes pratiques \textbf{REST} et en intégrant l'authentification \textbf{OIDC/Azure AD} avec contrôle d'accès basé sur les rôles.
                \item Effectué un \textbf{QA} de bout en bout sur les composants de la plateforme en exécutant des tests fonctionnels, en identifiant les défauts, en résolvant les bugs et les vulnérabilités de sécurité, et en améliorant la stabilité de l'application avant les mises en production.
            \end{highlights}
        \end{onecolentry}
        \vspace{0.2 cm}

        \begin{twocolentry}{
            Juin 2024 – Déc 2025
        }
        \textbf{Développeur Full Stack} -- \textbf{\hrefWithoutArrow{https://www.ariastechsolutions.com/}{Arias Tech Solutions}} \end{twocolentry}
        \vspace{0.10 cm}

        \begin{onecolentry}
            \begin{highlights}
                \item Développé des composants UI modulaires et des tableaux de bord réactifs avec \textbf{Angular}, \textbf{Tailwind CSS} et \textbf{PrimeNG}, permettant la visualisation de données low-code pour les utilisateurs métier.
                \item Construit des services backend évolutifs avec \textbf{NestJS} dans une architecture monorepo, intégrant des données provenant de \textbf{REST APIs}, \textbf{WebSockets} et bases de données \textbf{SQL/NoSQL}.
                \item Conçu des pipelines de données temps réel avec \textbf{Kafka} et \textbf{MQTT} pour traiter la télémétrie des capteurs océaniques afin de surveiller les écosystèmes marins.
                \item Développé des visualisations géospatiales dynamiques avec \textbf{Mapbox GL JS}, et implémenté des systèmes d'alerte basés sur des seuils avec notifications multicanaux.
                \item Contribué au développement de fonctionnalités d'analytique prédictive, de rapports automatisés et d'authentification et autorisation sécurisées avec \textbf{Keycloak}, supportant le contrôle d'accès basé sur les rôles sur les plateformes de données environnementales.
            \end{highlights}
        \end{onecolentry}
        \vspace{0.2 cm}


    % =========================== EDUCATION ==========================
    \section{Formation}
        \begin{twocolentry}{
            Oct 2022 – Déc 2025
        }
        \textbf{\href{https://um6p.ma/}{UM6P} - \href{https://1337.ma/en/}{1337}\href{https://www.42network.org/}{(42 Network)}} | Diplôme en informatique\end{twocolentry}
        \vspace{0.10 cm}

        Diplômé d'un programme intensif de génie logiciel axé sur les projets, portant sur le développement full-stack, les algorithmes, la conception de systèmes et la mise en réseau. Acquit une expérience pratique approfondie en programmation (C, C++, JavaScript), structures de données, bases de données, systèmes d'exploitation, protocoles réseau, sécurité et pratiques de développement logiciel collaboratif (Linux, Bash, Git et Docker). L'accent est mis sur la résolution de problèmes, l'apprentissage entre pairs et l'application concrète des principes d'ingénierie.

    % =========================== PROJECTS ==========================
    % \section{Projets}
    %     \begin{onecolentry}
    %         \begin{highlights}
    %         \item \textbf{\href{https://github.com/mtellami/ft_transcendence}{Application web de jeu Pong}}\\
    %         Un jeu Pong multijoueur full-stack entièrement développé en TypeScript, utilisant React.js et Tailwind CSS pour une interface réactive et moderne, et NestJS, PostgreSQL et Prisma pour le backend.\\Travaillé en équipe pour implémenter le gameplay en temps réel, la fonctionnalité de chat, les profils utilisateurs et l'authentification. L'accent est mis sur la conception modulaire, l'architecture propre et le développement d'API sécurisées tout en maintenant une collaboration fluide via Git et les workflows Agile.
    %         \vspace{0.10 cm}
    %
    %         \item \textbf{\href{https://github.com/mtellami/ft_transcendence}{Serveur HTTP}}\\
    %         Construit un serveur HTTP personnalisé en C++ dans le cadre d'un projet d'équipe, en appliquant les principes de programmation orientée objet (POO) et la programmation de sockets TCP de bas niveau pour gérer les requêtes et réponses HTTP. Compris en profondeur le fonctionnement interne des serveurs web, notamment la gestion des connexions, l'analyse, le routage et la génération de réponses, inspiré de NGINX.
    %         \end{highlights}
    %     \end{onecolentry}

    % =========================== SKILLS ==========================
    \section{Compétences}

        \begin{onecolentry}
            \textbf{Langages :} Java, Python, TypeScript, JavaScript, SQL, C/C++
        \end{onecolentry}
        \vspace{0.2 cm}
        
        \begin{onecolentry}
            \textbf{Frontend :} React, Next.js, Angular, Tailwind CSS, PrimeNG, Shadcn, TanStack
        \end{onecolentry}
        \vspace{0.2 cm}
        
        \begin{onecolentry}
            \textbf{Backend :} SpringBoot, Django, FastAPI, NestJS, Express.js, WebSockets, Kafka
        \end{onecolentry}
        \vspace{0.2 cm}
        
        \begin{onecolentry}
            \textbf{Bases de données :} PostgreSQL, MySQL, MongoDB, Redis
        \end{onecolentry}
        \vspace{0.2 cm}
        
        \begin{onecolentry}
            \textbf{Outils :} Git, Docker, Bash, Linux, Azure DevOps, Jira
        \end{onecolentry}
        \vspace{0.2 cm}

    % =========================== LANGUAGES ==========================
    \section{Langues}
        \textbf{Arabe :} Langue maternelle
        \hspace{1 cm}
        \textbf{Anglais :} Courant
        \hspace{1 cm}
        \textbf{Français :} Notions élémentaires

\end{document}
